\chapter{HER2/neu Tumor Marker}

\section*{What Is This Test?}
HER2 (human epidermal growth factor receptor 2), also known as HER2/neu or ERBB2, is a transmembrane receptor tyrosine kinase belonging to the epidermal growth factor receptor (EGFR/ErbB) family. The HER2 gene (ERBB2) is located on chromosome 17q12 and encodes a 185-kDa transmembrane glycoprotein with an extracellular ligand-binding domain, a transmembrane domain, and an intracellular tyrosine kinase domain. Unlike other ErbB family members, HER2 has no known direct ligand; instead, it functions as the preferred dimerization partner for other activated ErbB receptors (EGFR/HER1, HER3, HER4). HER2 amplification and/or protein overexpression occurs in approximately 15--20\% of invasive breast cancers and is associated with aggressive tumor biology, high histologic grade, increased proliferation rate, increased risk of metastases, and shortened overall survival in the absence of targeted therapy. HER2 testing is a mandatory component of the workup for all newly diagnosed invasive breast cancers (and recurrent/metastatic breast cancers) because HER2 positivity determines eligibility for HER2-directed therapies: trastuzumab (Herceptin), pertuzumab (Perjeta), ado-trastuzumab emtansine (T-DM1, Kadcyla), trastuzumab deruxtecan (T-DXd, Enhertu), lapatinib (Tykerb), neratinib (Nerlynx), and tucatinib (Tukysa). HER2 is also tested in gastric and gastroesophageal junction adenocarcinomas (15--25\% are HER2-positive) and, increasingly, in other tumor types (colorectal, lung, salivary gland, endometrial).

\section*{Alternative Names}
HER2; HER2/neu; ERBB2; Human Epidermal Growth Factor Receptor 2; c-erbB-2; HER2 Oncogene; HER2 Amplification; HER2 Overexpression; Immunohistochemistry for HER2 (IHC); Fluorescence In Situ Hybridization for HER2 (FISH)

\section*{Principle}
HER2 testing is performed directly on tumor tissue (breast biopsy or surgical resection specimen), not on blood (serum HER2 extracellular domain testing is available but not recommended for routine clinical use by ASCO/CAP). Two complementary methods are used, following the ASCO/CAP 2018 Guideline algorithm: immunohistochemistry (IHC) and in situ hybridization (ISH) \cite{wolff2018human}. IHC uses anti-HER2 antibodies (e.g., 4B5 clone, HercepTest) to stain formalin-fixed, paraffin-embedded tumor sections. A pathologist scores membrane staining intensity and completeness on a scale: 0 (negative, no staining or <10\% incomplete faint staining), 1+ (negative, >10\% incomplete faint staining), 2+ (equivocal, >10\% weak to moderate complete membrane staining, or strong complete staining in $\leq$10\%), and 3+ (positive, >10\% strong complete circumferential membrane staining). IHC 2+ cases reflex to ISH. ISH uses a fluorescent (FISH) or chromogenic (CISH, SISH) DNA probe for the HER2 gene locus (17q12) and a centromere enumeration probe for chromosome 17 (CEP17). The HER2/CEP17 ratio and the average HER2 copy number per cell are calculated. HER2/CEP17 ratio $\geq$2.0 AND average HER2 copy number $\geq$4.0 signals/cell = positive (amplified). Serum HER2 extracellular domain (ECD) can be measured by ELISA or chemiluminescent immunoassay, but this is not recommended for routine clinical decision-making.

\section*{History}
The HER2/neu oncogene was discovered in 1984 by Robert Weinberg's group at MIT and independently by the Genentech group. The gene was identified as a transforming oncogene in rat neuroblastoma (hence ``neu''). The human homolog (HER2/ERBB2) was cloned and characterized in subsequent years. Dennis Slamon and colleagues at UCLA demonstrated in 1987 that HER2 amplification in breast cancer correlated with shorter disease-free and overall survival. The development of trastuzumab, a humanized mouse monoclonal antibody targeting the HER2 extracellular domain, was a paradigm-shifting discovery by Genentech. Trastuzumab received FDA approval in 1998 for HER2-positive metastatic breast cancer \cite{slamon2001use}. The landmark adjuvant trials (NSABP B-31, NCCTG N9831, BCIRG 006, HERA) demonstrated that adding trastuzumab to adjuvant chemotherapy reduced the risk of breast cancer recurrence by 50\% and death by 33\% in HER2-positive early breast cancer, establishing HER2-directed therapy as one of the greatest advances in oncology. The ASCO/CAP HER2 Testing Guidelines were first published in 2007, updated in 2013, and most recently updated in 2018 to refine the equivocal ranges and align with outcomes from clinical trials.

\section*{Why This Test Is Ordered}
\begin{itemize}
  \item Determination of HER2 status in all newly diagnosed invasive breast cancers (and recurrences) to guide targeted therapy with HER2-directed agents (trastuzumab, pertuzumab, T-DM1, T-DXd, lapatinib, neratinib, tucatinib). HER2 testing is mandated by ASCO/CAP guidelines for every invasive breast cancer, regardless of tumor stage, histologic type, or hormone receptor status
  \item Mandatory repeat HER2 testing on a metastatic recurrence biopsy specimen if clinically feasible, because HER2 status can change between primary and metastatic tumors (HER2 discordance occurs in 5--15\% of cases)
  \item Determination of HER2 status in advanced gastric and gastroesophageal junction adenocarcinoma to identify candidates for trastuzumab-based chemotherapy (per the ToGA trial) \cite{bang2010trastuzumab}. HER2 positivity in gastric/GEJ cancer is defined by IHC 3+ or IHC 2+ with FISH-positive (HER2/CEP17 ratio $\geq$2.0); the criteria differ slightly from breast cancer
  \item Testing in other tumor types (colorectal cancer, non-small cell lung cancer, salivary gland carcinoma, endometrial serous carcinoma, urothelial carcinoma) in the context of comprehensive genomic profiling for clinical trials or off-label targeted therapy consideration, recognizing that HER2-directed therapies are not standard of care for most of these indications
  \item Repeat HER2 testing on residual disease after neoadjuvant chemotherapy if the original biopsy had low tumor cellularity or if testing was equivocal
\end{itemize}

\section*{Primary vs Secondary Test}
This is a primary biomarker test required for all patients with newly diagnosed invasive breast cancer. HER2 status is one of the three essential classifiers for breast cancer biology alongside estrogen receptor (ER) and progesterone receptor (PR). HER2 status determines eligibility for life-prolonging targeted therapy and is a component of the AJCC prognostic stage.

\section*{How This Test Is Performed}
HER2 testing is NOT performed on a blood sample. It is performed on tumor tissue obtained by core needle biopsy or surgical excision. The tissue is fixed in 10\% neutral buffered formalin for 6 to 72 hours, processed, and embedded in paraffin. A tissue section is cut and tested by IHC (HER2 protein expression). If the IHC result is equivocal (2+), the same tissue block is reflexively tested by FISH (HER2 gene amplification). In some centers, dual-probe FISH is performed as the primary HER2 test alongside or in place of IHC. Results are typically available in 3--7 days. The test requires a biopsy or surgical tissue specimen; it cannot be performed on a blood sample.

\section*{What Preparation Is Needed}
No patient preparation is required for the tissue biopsy procedure beyond standard pre-procedure instructions. The critical pre-analytical factor is adequate tissue fixation: the tissue must be fixed in 10\% neutral buffered formalin for at least 6 hours and no more than 72 hours. Under-fixation (and over-fixation) can produce false-negative and false-positive results. Tissue processing should follow standard CAP/CLIA protocols. Bone biopsies that undergo decalcification with strong acids can degrade HER2 protein and DNA, potentially producing false-negative results by both IHC and FISH; EDTA-based decalcification is preferred for metastatic bone biopsies when HER2 testing is planned.

\section*{Contraindications and Precautions}
\begin{itemize}
  \item The critical limitation of HER2 testing is pre-analytical tissue quality. False-negative results can occur due to: cold ischemic time >1 hour before fixation, under-fixation (<6 hours in formalin), poor tumor cell content in the biopsy, or sampling error (core needle biopsy of a heterogeneous tumor). False-positive results can occur due to: over-fixation (>72 hours in formalin), edge artifact or crush artifact on IHC interpretation, or polysomy of chromosome 17 (interpreted as increased HER2 copy number without true gene amplification). Equivocal results (IHC 2+) occur in approximately 15--20\% of cases and must be reflex tested by ISH. Genetic heterogeneity of HER2 amplification (HER2 amplification present in only a subset of tumor cells) can complicate interpretation. The serum HER2 extracellular domain (ECD) assay is NOT recommended for clinical decision-making by ASCO/CAP guidelines because its sensitivity and specificity are inadequate and it does not recapitulate tissue HER2 status. All laboratories performing HER2 testing must participate in proficiency testing programs (CAP HER2 proficiency testing) and must validate their assays per ASCO/CAP standards.
\end{itemize}
\section*{Risks}
Risks are related to the tissue biopsy procedure (core needle biopsy) rather than the laboratory test itself: pain, bruising, bleeding, infection at the biopsy site, and rarely pneumothorax (for breast lesions near the chest wall). The greater risk is a clinically significant false-negative or false-positive result that leads to withholding or administering HER2-directed therapy inappropriately.

\section*{What Happens After the Test}
Results are reported as HER2-positive (IHC 3+ or ISH amplified), HER2-negative (IHC 0 or 1+, or ISH non-amplified), or HER2-equivocal (if ISH result falls into an equivocal group per ASCO/CAP 2018 criteria). HER2-positive results should trigger multidisciplinary oncology evaluation for initiation of HER2-directed therapy (neoadjuvant, adjuvant, or metastatic). HER2-equivocal results require discussion with the pathologist and may warrant repeat testing on a different tissue block or re-biopsy. All HER2 test reports must include the specific assay performed (IHC antibody clone, ISH probe type), the scoring result, the interpretation (positive/negative/equivocal), and the laboratory's interpretation criteria (ASCO/CAP 2018).

\section*{Understanding Results}

\subsection*{IHC 0 or 1+ (HER2-Negative by IHC)}
\begin{itemize}
  \item IHC 0: No staining observed, or incomplete, faint/barely perceptible membrane staining in $\leq$10\% of tumor cells. HER2-negative. The patient is not eligible for HER2-directed therapy
  \item IHC 1+: Incomplete, faint/barely perceptible membrane staining in >10\% of invasive tumor cells. HER2-negative per ASCO/CAP 2018 guidelines. These patients are not eligible for conventional HER2-directed therapy (trastuzumab, pertuzumab, T-DM1). However, with the emergence of T-DXd (trastuzumab deruxtecan), IHC 1+ and IHC 2+/ISH-negative tumors are now classified as ``HER2-low,'' and the DESTINY-Breast04 trial demonstrated that T-DXd significantly improves progression-free and overall survival compared to physician's choice chemotherapy in HER2-low metastatic breast cancer \cite{modi2022trastuzumab}. Thus, IHC 1+ now has therapeutic implications in the metastatic setting
  \item ER-positive, PR-positive, and HER2-negative: Luminal A or Luminal B breast cancer (distinguished by proliferation index, Ki-67, and genomic assays). These patients are eligible for endocrine therapy (tamoxifen, aromatase inhibitors) with or without chemotherapy based on genomic risk profiling (Oncotype DX, MammaPrint)
\end{itemize}

\subsection*{IHC 2+ (Equivocal by IHC) --- Must Reflex to ISH}
\begin{itemize}
  \item Weak to moderate complete circumferential membrane staining in >10\% of invasive tumor cells. This result is NOT diagnostic; it requires reflex in situ hybridization (FISH) testing on the same tissue block. Approximately 15--20\% of IHC results are 2+, and of these, approximately 20--30\% are ISH-positive (truly HER2-amplified). The remaining 70--80\% are ISH-negative (HER2 non-amplified)
  \item If ISH is positive (HER2/CEP17 ratio $\geq$2.0 OR HER2 copy number $\geq$6.0 [group 1], or ratio $\geq$2.0 with copy number <4.0 [group 2], per ASCO/CAP 2018), the tumor is classified as HER2-positive
  \item If ISH is negative (HER2/CEP17 ratio <2.0 AND copy number <4.0 [group 5]), the tumor is classified as HER2-negative. These tumors are HER2-low by the Destiny-Breast04 criteria (IHC 1+ or IHC 2+/ISH-negative)
  \item Group 3 (ratio <2.0, copy number $\geq$6.0): HER2-positive if the concurrent IHC is 3+; HER2-negative if concurrent IHC is 2+
  \item Group 4 (ratio <2.0, copy number $\geq$4.0 and <6.0): HER2-negative if concurrent IHC is 2+; most of these represent chromosome 17 polysomy (multiple copies of chromosome 17) rather than true HER2 gene amplification
\end{itemize}

\subsection*{IHC 3+ (HER2-Positive) --- No ISH Needed}
\begin{itemize}
  \item Strong complete circumferential membrane staining in >10\% of invasive tumor cells. This is HER2-positive by definition per ASCO/CAP 2018; no reflex ISH testing is required (although some laboratories perform concurrent or reflex FISH for quality assurance). The patient is a candidate for HER2-directed therapy
  \item ISH positive (HER2/CEP17 ratio $\geq$2.0 AND average HER2 copy number $\geq$4.0 signals/cell): Gene amplification confirmed. HER2-positive breast cancer. These patients should receive neoadjuvant chemotherapy with dual HER2 blockade (trastuzumab + pertuzumab) followed by adjuvant therapy based on pathologic response at surgery (residual disease triggers T-DM1 per the KATHERINE trial \cite{vonminckwitz2019trastuzumab})
  \item HER2-positive breast cancer accounts for 15--20\% of all invasive breast cancers. It is more common in younger women, higher histologic grade tumors, and ductal histology. In the absence of HER2-directed therapy, HER2-positive breast cancer has an aggressive natural history with a median survival of <3 years in metastatic disease. With modern HER2-directed therapy, outcomes have been transformed: the 10-year disease-free survival for stage I--III HER2-positive breast cancer approaches 75--90\%
  \item HER2-positive gastric/gastroesophageal junction adenocarcinoma (IHC 3+ or IHC 2+/FISH-positive): Trastuzumab plus chemotherapy is the standard of care for the first-line treatment of metastatic HER2-positive gastric/GEJ cancer (ToGA trial). Unlike breast cancer, the ASCO/CAP gastric HER2 guidelines require IHC 3+ (or 2+ with FISH-positive) with a different scoring system adapted to the incomplete basolateral membrane staining typical of gastric cancers
\end{itemize}

\section*{Normal Results}
\begin{itemize}
  \item \textbf{Normal (HER2-negative) breast tissue:} IHC 0 or 1+. Normal breast epithelial cells express HER2 at a low physiologic level without gene amplification
  \item \textbf{HER2 IHC reporting (ASCO/CAP 2018 guideline):}
  \begin{itemize}
    \item IHC 0: No staining or <10\% incomplete faint membrane staining --- Negative
    \item IHC 1+: >10\% incomplete faint membrane staining --- Negative (HER2-low)
    \item IHC 2+: >10\% weak-moderate complete membrane staining --- Equivocal (reflex to FISH)
    \item IHC 3+: >10\% strong complete membrane staining --- Positive
  \end{itemize}
  \item \textbf{HER2 FISH reporting (ASCO/CAP 2018 guideline):}
  \begin{itemize}
    \item HER2/CEP17 ratio $\geq$2.0 AND HER2 copy number $\geq$4.0 (Group 1) --- Positive
    \item HER2/CEP17 ratio $\geq$2.0 AND HER2 copy number <4.0 (Group 2) --- In most cases, positive
    \item HER2/CEP17 ratio <2.0 AND HER2 copy number $\geq$6.0 (Group 3) --- Positive if concurrent IHC 3+
    \item HER2/CEP17 ratio <2.0 AND HER2 copy number $\geq$4.0 and <6.0 (Group 4) --- Negative
    \item HER2/CEP17 ratio <2.0 AND HER2 copy number <4.0 (Group 5) --- Negative
  \end{itemize}
  \item \textbf{Serum HER2 extracellular domain (ECD):} <15 ng/mL (>15 ng/mL is considered elevated for HER2-positive breast cancer monitoring but IS NOT recommended for routine clinical use by ASCO/CAP)
\end{itemize}

\section*{Follow-up Tests}
\begin{itemize}
  \item \textbf{Estrogen receptor (ER) and progesterone receptor (PR) IHC:} All breast cancers should be tested for ER and PR status simultaneously with HER2. Triple-negative breast cancer (ER-negative, PR-negative, HER2-negative) accounts for 10--15\% of breast cancers and is treated with chemotherapy and immunotherapy (pembrolizumab for PD-L1-positive TNBC). ER-positive and/or PR-positive disease is eligible for endocrine therapy
  \item \textbf{Ki-67 proliferation index:} Measured by IHC on the same tumor block. Provides information on tumor proliferation and is used to distinguish Luminal A (Ki-67 low, <14--20\%) from Luminal B (Ki-67 high) breast cancers
  \item \textbf{Genomic assays (Oncotype DX 21-gene recurrence score, MammaPrint 70-gene signature, Prosigna PAM50, EndoPredict):} Performed on ER-positive, HER2-negative breast cancers to estimate the risk of recurrence and guide decisions about the benefit of adding adjuvant chemotherapy to endocrine therapy
  \item \textbf{Comprehensive genomic profiling (tissue NGS, FoundationOne CDx or similar):} For metastatic breast cancer to detect other actionable targets beyond HER2: PIK3CA mutations (alpelisib eligibility), ESR1 mutations (acquired resistance to aromatase inhibitors), BRCA1/BRCA2 mutations (PARP inhibitor eligibility), NTRK fusions, MSI-H/dMMR (pembrolizumab eligibility), and TMB-High
  \item \textbf{Germline BRCA1 and BRCA2 genetic testing:} All patients with HER2-negative metastatic breast cancer (particularly ER-positive/HER2-negative and triple-negative) should be offered germline BRCA testing for PARP inhibitor (olaparib, talazoparib) eligibility. Patients with HER2-positive breast cancer who have a strong family history should also be offered testing
  \item \textbf{Repeat HER2 testing on metastatic recurrence biopsy:} At the time of disease progression to evaluate for loss of HER2 expression (which can occur in 10--15\% of patients and indicates resistance to ongoing HER2-directed therapy) or gain of HER2 expression (which would make patients eligible for HER2-directed therapy). Discordance in HER2 status between primary and metastatic disease occurs in 5--15\%
  \item \textbf{Brain MRI:} HER2-positive breast cancer has a high propensity for brain metastases (30--50\% lifetime incidence in metastatic HER2-positive disease). CNS imaging should be performed at the time of metastatic presentation
  \item \textbf{Cardiac monitoring (transthoracic echocardiogram or MUGA scan):} Baseline left ventricular ejection fraction (LVEF) assessment before initiating HER2-directed therapy. Serial monitoring every 3--4 months during trastuzumab-based therapy to detect cardiotoxicity (trastuzumab-induced cardiomyopathy, an asymptomatic decline in LVEF $\geq$10 percentage points to <50\%)
\end{itemize}

\section*{References}
\bibliographystyle{unsrtnat}
\bibliography{references}

\clearpage
\section*{Regulatory Status and Authorization Date}
This chapter describes a test category, not a specific commercial product. FDA, CDSCO, EU IVDR, and MHRA approval or registration dates therefore cannot be assigned without the assay manufacturer, product name, instrument, and intended use. For laboratory-developed tests, an individual agency approval date may not exist. See the Preface for the interpretation of regulatory status.

\clearpage
